% !PRINTED 424
\weberpage{424}

%<<CONT
und das Vorzeichen hiervon ist von $k$ abhängig. Nehmen wir aber $k = 1$, also $r = r_0$, so sind alle die Winkel
\[
\frac{2\pi}{n}, \frac{4\pi}{n} \ldots \frac{(n-1)\pi}{n}
\]
zwischen Null und $\pi$ gelegen und die Sinus alle positiv. Wir schliessen hieraus, dass in diesem Falle
\weq{14}{P = (-i)^{\frac{n-1}{2}} n^{\frac{n-2}{2}} \sqrt{n}}
ist, und $\sqrt{n}$ positiv zu nehmen ist.

Das Vorzeichen des Productes der $R$ hängt von der Wahl von $r$ nicht ab, sondern nur von der Beschaffenheit von $n$. Es ist daher von Interesse, das Product $Q$, dessen Vorzeichen von der Wahl von $r$ abhängt, für sich zu bestimmen. Wir erhalten es auf folgende Weise:

Multipliciren wir den Ausdruck (11) mit
\[
r \cdot r^2 \cdot r^3 \cdots r^{\frac{n-1}{2}} = r^{\frac{n^2-1}{8}}
\]
und sodann mit
\[
r^{-1} \cdot r^{-2} \cdot r^{-3} \cdots r^{-\frac{n-1}{2}} = r^{-\frac{n^2-1}{8}},
\]
so erhalten wir
\[
r^{\frac{n^2-1}{8}} Q = (r^2 - 1)(r^4 - 1) \cdots (r^{n-1} - 1)
\]
\[
r^{-\frac{n^2-1}{8}} Q = (1 - r^{n-2})(1 - r^{n-4}) \cdots (1 - r).
\]
Da nun die Exponenten $2, 4 \ldots n-1, n-2, n-4 \ldots 1$ zusammen alle Zahlen $1, 2 \ldots n-1$ umfassen, so ergiebt sich durch Multiplication dieser beiden Ausdrücke:
\[
Q^2 = (-1)^{\frac{n-1}{2}} \Pi(r-1) = (-1)^{\frac{n-1}{2}} n,
\]
also
\weq{15}{Q = \pm i^{\frac{n-1}{2}} \sqrt{n}.}
Die Vergleichung mit (13) ergiebt
\weq{16}{2^{\frac{n-1}{2}} \sin\frac{2\pi k}{n} \sin\frac{4\pi k}{n} \cdots \sin\frac{(n-1)\pi k}{n} = \pm \sqrt{n}.}
Das Vorzeichen in dieser Formel hängt, wie in (15), noch von $k$ ab; es ist aber das positive, wenn $k = 1$ ist. Im Uebrigen wollen wir über dies Vorzeichen, das weiterhin noch genauer untersucht werden wird, noch einen wichtigen Satz ableiten.
